\documentclass[../principal.tex]{subfiles}
\graphicspath{{\subfix{../imagenes/}}}

\begin{document}

\chapter{Filtros}

Los filtros son máquinas que operan a nivel de la frecuencia de una señal.

Existen muchos tipos de filtros, según cómo operan en la frecuencia:

\begin{itemize}
    \item Pasabajos
    \item Pasaaltos
    \item Pasabanda
    \item Rechazabanda
\end{itemize}

A nivel de su funcionamiento y alimentación, existen filtros que solamente son capaces de atenuar la señal, y otros que también pueden amplificarla. Reciben los nombres de filtros pasivos y activos respectivamente.

\section{VCF}

Existe en el rubro de máquinas sonoras un tipo de filtro controlado por voltaje, que recibe el nombre VCF por sus siglas en inglés.

\section{Filtro pasabajos}

Un filtro pasabajos atenúa la señal según su contenido en frecuencia, permitiendo que las componentes que están bajo una cierta frecuencia de corte no tengan cambios en su amplitd, mientras que las que estén por sobre sean atenuadas.

El filtro pasabajos es un clásico de las máquinas sonoras ya que emula el comportamiento del sonido cuando se propaga en el espacio, donde los sonidos graves son capaces de recorrer grandes distancias sin disipaciones considerables, mientras que las frecuencias agudas se atenuan con mayor dramatismo.

\newpage

\end{document}

